\begin{table}[htbp]
\centering
\small
\caption{两组与三组方案的工作量比较}
\label{tab:q4-workload}
\setlength{\tabcolsep}{6pt}
\renewcommand{\arraystretch}{1.2}
\begin{tabular}{llrrrr}
\toprule
分组规模 & 任务组 & 货箱/箱 & 联合架次 & 能耗/kWh & 任务占用时长/h \\
\midrule
2组 & G1 & 48 & 13 & 40.68 & 8.09 \\
2组 & G2 & 32 & 12 & 28.28 & 6.94 \\
3组 & G1 & 22 & 8 & 19.71 & 4.55 \\
3组 & G2 & 29 & 10 & 29.41 & 6.46 \\
3组 & G3 & 29 & 7 & 19.85 & 4.02 \\
\bottomrule
\end{tabular}
\par\smallskip
\begin{minipage}{0.98\linewidth}
\footnotesize 注：采用货箱数和联合架次数偏离组均值不超过25\%时的组内重派方案。任务占用时长为各运输任务开始至返航、中继任务准备开始至返航时长之和，允许任务并行；它不等于联合任务完成时间。
\end{minipage}
\end{table}
